\section{Mechanism and accounting definitions}
\label{sec:mechanism}
\begin{definition}[World and observations]
A world consists of a registry, pending registrations, pending terminations and
a cumulative treasury counter. Its reachability starts from any registry
satisfying the original registry reachability relation, with both pending pools
empty and the counter zero. It then closes under accepted bookings, rejections,
retractions and economic folds. It does not assert that every initial registry
was produced by the economic transactions of this paper.
\end{definition}
An observed economic transaction pairs ledger-shaped inputs, references,
outputs, fee and donation with request contents and typed token movements.
The observation checks bind these fields to an unspent set and the world.
This is a logical observation relation, not a decoder for a Cardano transaction.
The definition is \decl{Singular.WorldReachable}; the accepted judgements are
\decl{Singular.book}, \decl{Singular.bookTermination},
\decl{Singular.rejectionAccepted}, \decl{Singular.retractionAccepted} and
\decl{Singular.foldAccepted}, at the pin above. These are definitions, not
axiom-checked theorem claims; their consequences are cited below.

\begin{definition}[Amounts and clocks]
Let $R$ be the registration charge, $t_{\min}$ the minimum tip, $\delta$ the
maximum booking width, and $W_p,W_r$ the processing and reclaim windows.
These and the registry identity and policy bindings are fixed at creation.
The minimum tip occupies the existing wire key \texttt{maxFee}; it is no
longer an upper fee pin. A request declares tip $t\ge t_{\min}$ and carrier
amount $a$. Its held lovelace is $H$. A registration is booked in a nonempty
interval $[L,U)$ with $U-L\le\delta$, and records $s=U$.
\end{definition}
The carrier amount is declared under the ledger's output-minimum rules; it
is not an additional registry parameter. The model checks exact declared
amounts, not the actual minimum or the fee needed by a concrete transaction.
The original application token's carrier and a later termination request's
carrier need not have equal lovelace amounts: their output shapes differ.

\begin{table}[ht]
\caption{Registration request. Extra assets or hidden lovelace surplus are refused.}
\centering\small
\begin{tabular}{p{.25\linewidth}p{.65\linewidth}}\toprule
Field & Required contents\\\midrule
Lovelace & Exactly $H=R+t+a$; $t\ge t_{\min}$.\\
Payload & Register, key, refund address and destination.\\
Datum & Separately booked destination datum; copied by the fold, not included in these tag-name commitments.\\
Application tag & Exactly one, bound to the payload and nonce; minted in this booking.\\
Trie-ready tag & Exactly one, bound to the operation and the same nonce; minted in this booking.\\
Clock & $s=U$, $L<U$, $U-L\le\delta$.\\
Exits & Fold; signed retraction; reference-read rejection.\\\bottomrule
\end{tabular}
\end{table}
\begin{table}[ht]
\caption{Termination request. Application consent is release from custody.}
\centering\small
\begin{tabular}{p{.25\linewidth}p{.65\linewidth}}\toprule
Field & Required contents\\\midrule
Lovelace & Exactly $H=t+a$; $t\ge t_{\min}$.\\
Singular tag & Exactly the genuine tag of this active key, moved from application custody.\\
Trie-ready tag & Exactly one, bound to its operation and nonce; minted in this booking.\\
Application tag & None.\\
Refund and clock & Neither a refund address nor a booked time.\\
Exit & Fold only, without a registration window.\\\bottomrule
\end{tabular}
\end{table}

\checked The request shapes, same-booking tag minting and bounded booking clock
are consequences of accepted booking and reachable pending membership.
\evidence{booked_request_has_exact_contents}
\evidence{booked_request_tags_commit_to_own_fields}
\evidence{booked_time_is_interval_end}
\evidence{booking_width_within_limit}
\evidence{termination_booking_has_exact_contents}
\evidence{booked_termination_tag_commits_to_own_operation}
Tag names equal the computed commitments; this is not a proof of collision
resistance. Application tags minted outside a booking are outside Singular's
supply guarantees and cannot replace the simultaneous mint required by booking.
The singular tag is a different asset: registration mints it at the destination,
termination booking moves it into a request, and termination folding burns it.
The state token identifies the registry output. The abstract identity spaces
exclude its equality with a trie-ready tag even if a commitment is zero.
\evidence{trie_ready_tag_is_not_state_token}

\subsection{Admitting an exit}
A registration fold's entire nonempty validity interval must lie in
$[s,s+W_p)$. A retraction's must lie in $[s+W_p,s+W_p+W_r)$ and its refund
address must sign. A rejection needs a current lifecycle mismatch, or a
nonempty interval beginning at or after $s+W_p+W_r$. Thus mismatch remains
available during either window, while a compatible request is protected
through both. These are interval-containment rules, not tests of a wall clock
supplied by the processor. Equality of an interval's upper bound with the
window endpoint is allowed because the interval is half-open.

\begin{center}\small
\begin{tabular}{llll}\toprule
Registration time & Fold & Retract & Reject\\\midrule
$[s,s+W_p)$ & Yes & No & Mismatch\\
$[s+W_p,s+W_p+W_r)$ & No & Signed & Mismatch\\
$[s+W_p+W_r,\infty)$ & No & No & Expiry or mismatch\\\bottomrule
\end{tabular}
\end{center}
The table describes admission windows, subject to all other checks. A funded
retraction exists in the model under the explicit premises stated below;
neither a window nor that existence proof guarantees ledger inclusion.

A retraction spends one pending registration and may spend ordinary lovelace
outputs to pay its fee. Extra inputs hold none of the four registry assets,
have neither the request nor state role, and carry no request datum. Spending
inputs are present, distinct and disjoint from references. The state is neither
read nor spent. The refund address signs and receives at least all held
lovelace, including the tip; both request tags burn. No tip is taken and
nothing goes to the treasury. The checked payment bound below states precisely
what ordinary change at another address may contain.

A rejection references the actual unspent state output: one state token with
the exact policy and name, state role, inline datum, complete matching
configuration and current root. It never spends that output. Its distinct
request inputs are exactly its members; every member is checked against that
same state and interval. One invalid member refuses the entire batch. A fold
that consumes the referenced output first invalidates this proposed rejection;
rebuilding can starve indefinitely. A rejected proposal that never enters the
ledger has no ledger transaction fee, though preparation has an external cost.
The state and protection results are stated in Section~\ref{sec:properties}.

The request validator governs spending; a tag policy's permissive burn branch
does not confer a spending right. Registration exits burn both request tags;
a termination fold burns its trie-ready and singular tags. An untagged output
at the request address is not a pending request and can be swept. The checked
sweep judgement preserves the world; a concrete script implementing that
distinction remains a correspondence obligation.
\evidence{sweep_leaves_pool_untouched}

\subsection{Allocation is not net payoff}
\begin{definition}[Fold allocations]
For a fold, let $I,O$ be the sums of ordinary input and output lovelace,
respectively, $f$ its ledger fee, $T=\sum t$ its tips, $n$ its registrations,
and $D$ its donation. State input and continuation values are equal.
Ordinary funding and change are lovelace only; an unrelated native asset is
refused. Carrier outputs hold exactly their declared amount, booked datum and newly
minted singular tag. Ordinary change is a separate output even if its address
equals a carrier's address.
\end{definition}
\checked
\begin{theorem}[Economic fold allocation]
Every accepted economic fold in a reachable world satisfies
\begin{align}
D &= nR+\sum_{q\text{ termination}}(H_q-t_q),\label{eq:donation}\\
O+f &= I+T,\qquad f\le I.\label{eq:gross}
\end{align}
It is balanced; all request tags are burned and the treasury counter increases
by $D$.
\end{theorem}
The mechanized proof is the following checked declaration. The table separates
request allocation from fees.
\evidence{fold_routes_every_amount}

\begin{table}[ht]
\caption{Request allocation. Rejection alone supplies no full balance proof.}
\centering\small
\begin{tabular}{p{.24\linewidth}p{.64\linewidth}}\toprule
Exit & Allocation of held lovelace and token fate\\\midrule
Registration fold & $R$ donation; $a$ carrier at destination; $t$ gross ordinary allocation. Burn both request tags; mint singular tag.\\
Registration rejection & Refund address receives at least $R+a$. The judgement proves no exact processor reward. Burn both request tags.\\
Registration retraction & At least $H$ refunded, including $t$; ordinary funding pays $f$; zero donation, no tip taken. Burn both request tags. Balanced.\\
Termination fold & $H-t=a$ donation; $t$ gross ordinary allocation. Burn trie-ready and singular tags; no refund or destination output.\\\bottomrule
\end{tabular}
\end{table}
Rejection and retraction refund evidence is given under Property~6 below.
The rejection promise is a refund floor, not an exact split. Its judgement
checks a lower bound and does not require full monetary
balance; neither an exact wage nor a rejection analogue of~\eqref{eq:gross}
is supplied. A processor may leave part of its tip to the requester;
over-refunding is compatible with that lower bound.

\checked For a retraction, let $I$ be its ordinary funding, $F$ all output
lovelace at its refund address and $P$ all output lovelace elsewhere. Every
accepted retraction in a reachable world satisfies
\[
 F\ge H,\qquad f\le I,\qquad D=0,\qquad P\le I-f,
\]
and is balanced with an unchanged treasury. Ordinary funding and every output
are lovelace only. The model has no wallet-ownership or output-purpose
observation: $P$ may be change or a payment to a third party, but neither can
consume the request's cost or tip. Saying that the requester pays assumes it
owns the ordinary funding; the checked theorem itself does not assign owners.
\evidence{retraction_pays_only_the_refund_and_the_fee}
\hand Conservation gives $F+P+f=H+I$. Subtracting $F\ge H$ gives the
outside-address bound; any over-refund consumes ordinary funding as well.

\hand For a rejection whose inputs are exactly its members, additionally
assume \decl{Singular.balanced} and zero donation. Partition its outputs into
those paid to the set of refund addresses, totaling $F$, and all other outputs,
totaling $P$. Distinct members and summed floors give $F\ge\sum H-T$.
Balance gives $F+P+f=\sum H$, hence $P+f\le T$, and therefore $P\le T$.
This bounds value outside refund addresses, not the wealth of a party that
also controls a refund address. Balance is an additional ledger fact here,
not a checked conclusion of rejection admission.

\hand With signed differences,~\eqref{eq:gross} gives $O-I=T-f$.
This follows by subtracting $I+f$ in the integers, not by truncated natural
subtraction. It is the net change of the ordinary funding pool. If a processor
owns that pool, pays the fee and has no other role, its monetary net is $T-f$;
including external cost gives $T-f-e$. An address shared with a carrier does
not make the carrier's value a processor reward.

For one request, assume separate funder, destination, processor and treasury;
write $b$ for the booking fee. Measure transfers from before booking, treat
assets and application services separately, and count carrier value as an
asset of the destination. The following illustrative accounting sums to
minus the ledger fees, before external costs. In the funded retraction row,
the funder owns the refund address and ordinary funding, and all remaining
funding returns to it. This is an accepted model construction under the
availability theorem's window, signing and funding premises, not every
possible accepted retraction.
\begin{center}\small
\begin{tabular}{lrrrr}\toprule
Outcome & Funder & Destination & Processor & Treasury\\\midrule
Registration fold & $-R-t-a-b$ & $a$ & $t-f$ & $R$\\
Termination fold & $-t-a-b$ & $0$ & $t-f$ & $a$\\
Rejection, exact refund\! & $-t-b$ & $0$ & $t-f$ & $0$\\
Funded retraction & $-b-f$ & $0$ & $0$ & $0$\\\bottomrule
\end{tabular}
\end{center}
\hand The first two rows sum to $-b-f$ by cancellation. The third does too
\emph{if} an additionally balanced rejection pays exactly $H-t$ to the
funder's refund address, donates zero and returns the residue to its processor.
These added hypotheses are not proved by the rejection judgement. Native
asset values and a preexisting singular tag's valuation are excluded.
The retraction row returns $H+I-f$ to the funder from its total $H+I+b$
outlay, giving $-b-f$; it pays no processor or treasury allocation.
\evidence{pending_registration_can_be_retracted}

For example, with $R=2{,}000{,}000$, $t=400{,}000$, $a=1{,}500{,}000$,
$b=200{,}000$ and $f=170{,}000$, a registration holds $3{,}900{,}000$.
Its four net entries are $-4{,}100{,}000$, $1{,}500{,}000$, $230{,}000$
and $2{,}000{,}000$, summing to $-370{,}000$. These are arithmetic examples,
not measured output minima, costs or recommended parameters.

\subsection{Donation encoding}
The model represents $D$ by a natural number; a zero donation is encoded by
omitting the optional donation field, not by serializing a zero value.
\evidence{zero_donation_omits_field}
Conway transaction-body field 22 is a positive coin donation; optional field
21 is the separate current-treasury-value assertion~\cite{ConwayCDDL}.
When the latter is present, the ledger checks it against current treasury
state~\cite{ConwayLedger}. Donation alone does not make that assertion or
introduce its stale-value hazard. This describes the cited ledger revision;
translation to a concrete transaction still needs correspondence evidence.
